\subsection{27.02.2026}

\subsubsection{Ziel}

Ziel des Tages war es, den JOY-PI mit dem Internet zu verbinden,
um anschließend die neuen Funktionen der Taupunktberechnung
testen zu können.

\subsubsection{Durchgeführte Arbeiten}

Zunächst wurde überprüft, ob die bestehenden Skripte auf dem
JOY-PI weiterhin fehlerfrei ausgeführt werden können.
Dabei zeigte sich, dass die bisherigen Funktionen weiterhin
ohne Fehlermeldungen funktionieren.

Anschließend wurde versucht, eine Internetverbindung für den
JOY-PI herzustellen, um auf das GitHub-Repository zugreifen
und die aktuellen Änderungen herunterladen zu können.

Da jedoch keine funktionierende Internetverbindung verfügbar
war, konnte keine Verbindung zum GitHub-Repository hergestellt
werden. Aus diesem Grund war es nicht möglich, die neu
entwickelten Funktionen der Taupunktberechnung auf dem
JOY-PI zu testen.

\subsubsection{Nächste Schritte}

In der nächsten Woche soll der zweite DHT11-Sensor
angeschlossen und in die Taupunktberechnung integriert
werden.

Außerdem sollen die Berechnungsergebnisse in die
Hauptlogik übernommen und anschließend auf dem
LCD-Display ausgegeben werden.